\emph{Setting.} Throughout, $(M,G)$ satisfies (T1) to (T3) of \claimref{sheaf.tameness}, with the definitions of \claimref{synthesis.wall-rung-count-is-height-of-dying-line-class}: a $k$-rung completion is a smooth symmetric form $h$ on $T_k = T^*M \oplus \underline{\R}^{2k}$ over $\Mp$, nondegenerate on $\Mp$ including $\Sig$, with $h|_{T^*M} = G$ along $\Sig$ and $h = G \perp H_k$ outside a collar. $N(\cdot)$ denotes a maximal negative subbundle and $\overline G$ the induced form on $T^*M|_\Sig/\shL$ (\claimref{sheaf.degeneration-sheaf}). Line bundles with equal $w_1$ are isomorphic and $w$ is multiplicative (Fact \claimref{sheaf.imported-real-line-bundles}).

\emph{Inputs.} (K1) is the hypothesis (L1) of the statement, cited from memory and not in the store. (K2) is (L2), which is now clause (6) of Theorem \claimref{synthesis.cohomology-of-the-quarter-turn-lens-spaces-from-the-equivariant-cell-structure}, proved there given the imported Atiyah-Hirzebruch filtration; the route recorded below is the literature's, kept for reference and not used. (K3) was the hypothesis (L3) of version 1 and is no longer used: $x^{2} = 0$ is clause (3) of the same theorem. That theorem is stated for the quarter turn $J$ in every coordinate plane; any orthogonal complex structure $g$ on $\R^{2n+2}$ is conjugate to $J$ by an orthogonal map (an orthonormal basis $u_{1}, g u_{1}, \dots, u_{n+1}, g u_{n+1}$), so $S^{2n+1}/\langle g \rangle$ is isometric to its $\Sig_{n}$, with the flat lines corresponding. (K1) Complex K-theory $K(X)$ of a finite complex is a commutative ring with lambda-operations: $\lambda_t(a) = \sum_{j} \lambda^{j}(a) t^{j}$ satisfies $\lambda_t(a + b) = \lambda_t(a)\lambda_t(b)$, $\lambda_t(\ell) = 1 + \ell t$ for a line bundle $\ell$, $\lambda_t(1) = 1 + t$, and $\lambda^{j}(E) = 0$ for $j > \operatorname{rank} E$ when $E$ is an actual bundle; the reduced ring $\widetilde K(X)$ is nilpotent (Atiyah, \emph{K-Theory}, 1967). (K2) For a finite group $\Gamma$ acting freely and linearly on the unit sphere $S(V)$ of a complex representation $V$, $K(S(V)/\Gamma) \cong R(\Gamma)/(\lambda_{-1}(V))$: the Gysin sequence of the equivariant Thom isomorphism $K_\Gamma(D(V), S(V)) \cong R(\Gamma)$, whose Thom class restricts to $\lambda_{-1}(V) = \sum_i (-1)^{i}\Lambda^{i}V$ on the zero section, together with $K^{1}_\Gamma(\mathrm{pt}) = 0$ and $K_\Gamma(S(V)) = K(S(V)/\Gamma)$ for a free action (Atiyah, ibid., Section 2.7; Atiyah 1961; Kambe 1966). For $\Gamma = \Z_{4}$ with character $\chi$ (the generator acting by a complex structure) and $V = (n+1)\chi$ this gives $K(L^{2n+1}(4)) = \Z[\chi]/(\chi^{4} - 1, (\chi - 1)^{n+1})$, the class $\chi^{e}$ being the flat line bundle on which the generator acts by $\chi^{e}$. (K3) The mod-2 reduction of a class in $H^{1}(X; \Z_{4})$ has $Sq^{1} = 0$, since $Sq^{1}$ is the Bockstein of $0 \to \Z_{2} \to \Z_{4} \to \Z_{2} \to 0$; for a degree-one class $Sq^{1}x = x^{2}$ (Hatcher, \emph{Algebraic Topology}, Section 4.L).

\emph{(1).} Let $h$ be a $k$-rung completion. By clause (1) of \claimref{synthesis.wall-rung-count-is-height-of-dying-line-class} there is a real bundle $N$ of rank $k - 1$ over $\Sig$ with $N(\overline G) \oplus \shL \oplus N \cong N(\overline G) \oplus \underline{\R}^{k}$. Complexify: writing $c(E)$ for the complexification of a real bundle $E$ (the bundle $E \oplus E$ with the complex structure $(y_1, y_2) \mapsto (-y_2, y_1)$, the unit bivector of each plane $\operatorname{span}(y, y')$ of a copy of $E \oplus E$), complexification commutes with direct sums, so in the group $K(\Sig)$, after cancelling $c(N(\overline G))$, $[c(N)] = k - \xi$ with $\xi := [c(\shL)]$. Since $\shL \otimes \shL$ is trivial, $\xi^{2} = 1$. By (K1),
\[
\lambda_t(c(N)) = \lambda_t(k)\,\lambda_t(\xi)^{-1} = (1+t)^{k}(1 + \xi t)^{-1} = (1+t)^{k} \sum_{j \geq 0} (-\xi)^{j} t^{j},
\]
the inverse being a formal power series. The coefficient of $t^{k}$ is $\sum_{j=0}^{k} \binom{k}{k-j}(-\xi)^{j} = (1 - \xi)^{k}$, and it vanishes because $c(N)$ is an actual bundle of complex rank $k - 1 < k$. From $\xi^{2} = 1$, $(1-\xi)^{2} = 2(1 - \xi)$, so $(1-\xi)^{k} = 2^{k-1}(1-\xi)$; hence $2^{k-1}(1-\xi) = 0$. If $1 - \xi$ has order $2^{a}$ in $K(\Sig)$, then $k - 1 \geq a$. Nothing in this argument depends on $q$: the obstruction is the same on definite and indefinite sides. (The check computes the coefficient of $t^{k}$ in $\Z[\xi]/(\xi^{2} - 1)$ and finds $2^{k-1}(1-\xi)$ for $k = 1, \dots, 5$.)

\emph{(2).} Let $\Sig = L^{2n+1}(4) = S^{2n+1}/\langle g \rangle$, $g$ acting by a complex structure on $\R^{2n+2}$. Then $\pi_1(\Sig) = \Z_{4}$, $H^{1}(\Sig; \Z_{2}) = \operatorname{Hom}(\Z_{4}, \Z_{2}) = \Z_{2}$, and the line bundle $\shL$ with $x \neq 0$ is the flat line on which $g$ acts by $-1$. Its class $x$ is the reduction of the generator of $H^{1}(\Sig; \Z_{4}) = \operatorname{Hom}(\Z_{4}, \Z_{4})$, and $x^{2} = 0$ by clause (3) of Theorem \claimref{synthesis.cohomology-of-the-quarter-turn-lens-spaces-from-the-equivariant-cell-structure} (the cup coefficient of two odd cochains is $6$, even; in version 1 this step was $x^{2} = Sq^{1}x = 0$ by (K3)): $\operatorname{ht}(x) = 2$ for every $n$. The complexification of $\shL$ is the flat complex line on which $g$ acts by $-1 = \chi(g)^{2}$, so $\xi = \chi^{2}$ in (K2), and $1 - \xi = 1 - \chi^{2}$.

\emph{Lower bound for every $n$.} Let $I := e_1 e_2$ in $\Cl(2,0)$, so $I^{2} = -1$ and $\Z[I] = \{a + bI\}$ is the ring of Gaussian integers written with the GA unit; put $\pi := I - 1$. Then $\pi(-1 - I) = 2$ and $I\pi^{2} = I(-2I) = 2$, so $2 = I \pi^{2}$ with $I$ a unit. The assignment $\chi \mapsto I$ kills $\chi^{4} - 1$ ($I^{4} = 1$) and sends $(\chi - 1)^{n+1}$ to $\pi^{n+1}$, so it defines a ring map $\phi : K(\Sig) \to \Z[I]/(\pi^{n+1})$. It sends $2^{j-1}(1 - \chi^{2})$ to $2^{j-1}(1 - I^{2}) = 2^{j} = I^{j}\pi^{2j}$, whose $\pi$-adic valuation is exactly $2j$ ($\pi$ is prime of norm $2$; $a + bI$ is divisible by $\pi$ iff $a + b$ is even). So $2^{j-1}(1 - \xi) \neq 0$ whenever $2j < n + 1$. A $k$-rung completion needs $2^{k-1}(1-\xi) = 0$ by (1), hence $2k \geq n + 1$, that is $k \geq \lceil (n+1)/2 \rceil = 1 + \lfloor n/2 \rfloor$.

\emph{Exact orders for $n \leq 8$.} In the coordinates $s = \chi - 1$, $K(\Sig) = \Z[s]/(s^{n+1}, (1+s)^{4} - 1)$, and $\widetilde K(\Sig)$ is the quotient of $\Z s \oplus \dots \oplus \Z s^{n}$ by the lattice spanned by $s^{i}((1+s)^{4} - 1)$, $0 \leq i \leq n - 1$, truncated at $s^{n}$. This basis is triangular with diagonal entries $4$, so the lattice has index $4^{n}$, the order of $\widetilde K(L^{2n+1}(4))$ (clause (6) of the cited theorem: its Atiyah-Hirzebruch filtration quotients are $H^{2j}(\Sig; \Z) = \Z_{4}$, $1 \leq j \leq n$), and membership is decided exactly by back substitution. With $1 - \chi^{2} = -(2s + s^{2})$, the least $a$ with $2^{a}(1 - \chi^{2})$ in the lattice is $1, 1, 2, 2, 3, 3, 4, 4$ for $n = 1, \dots, 8$, that is $\lfloor (n+1)/2 \rfloor$. With (1), $k \geq 1 + \lfloor (n+1)/2 \rfloor$ for $n \leq 8$. In particular at $n = 3$: $2(1 - \chi^{2}) = -4s - 2s^{2}$ is not in the lattice spanned by $(4,6,4)$, $(0,4,6)$, $(0,0,4)$ (the coefficient of $s^{3}$ would need $6 + 4c = 4$), so $k(L^{7}(4)) \geq 3$. At $n = 1$, $2(1 - \chi^{2}) = -4s = -((1+s)^{4} - 1)$ modulo $s^{2}$, so the bound is $2$, as it must be: the $L(4,1)$ wall with an indefinite side has a two-rung completion (\claimref{synthesis.lens-space-wall-rung-count-sees-the-integral-twist}(1)). The same code returns the order $2^{n}$ of $\chi - 1$ in $\widetilde K(\R P^{2n+1})$, the classical value, as a control.

\emph{(3), the walls.} Let $n = 2r - 1$ and identify $\R^{2n+2} = \R^{4r}$ with $r$-tuples $q = (q_1, \dots, q_r)$ of even elements of $\Cl(3,0)$ (quaternions), with the inner product $\langle q, q' \rangle = \sum_a \langle q_a \tilde q'_a \rangle_0$. Let $B := e_2e_3$ and $C := e_3e_1$; then $B^{2} = C^{2} = -1$ and $BC = -CB$. Left multiplication $q \mapsto Bq$ (componentwise) is an isometry with square $-1$, a complex structure, so it generates a free $\Z_{4}$ action on the unit sphere with quotient $\Sig = L^{2n+1}(4)$. Put $f(q) := Cq$. Then $\langle q, f(q) \rangle = \sum_a \langle q_a \tilde q_a \tilde C \rangle_0 = |q|^{2}\langle \tilde C \rangle_0 = 0$ and $|f(q)| = |q|$, so $f$ is a unit tangent field on the sphere, and $f(Bq) = CBq = -B f(q)$. The differential of $q \mapsto Bq$ is again left multiplication by $B$, so $f$ is not a vector field on $\Sig$ but a nowhere-zero section of $\shL \otimes T\Sig$ for the line $\shL$ on which the generator acts by $-1$: its span is a line subbundle of $T\Sig \cong T^*\Sig$ (via the round metric $A$, which $B$ preserves) isomorphic to $\shL$, with $x \neq 0$. On $M = \Sig \times \R$ put
\[
G_t := \bigl(A - (1-t)\, A(\pi_f\,\cdot, \pi_f\,\cdot)\bigr) \perp (\epsilon\, dt^{2}), \qquad \epsilon = \pm 1,
\]
with $\pi_f$ the orthogonal projection on the span of $f$. It is single valued on $\Sig$, since $f$ enters quadratically. At $t = 0$ it has corank exactly one with radical the span of $f$, a smooth line bundle isomorphic to $\shL$; $G_t(f, f) = t|f|^{2}$, so the degeneration is first order and (T3) holds; the wall $\Sig \times \{0\}$ is a closed hypersurface separating $M$. For $t > 0$ the physical side has signature $(2n+1, 1)$ if $\epsilon = -1$ and $(2n+2, 0)$ if $\epsilon = +1$. The check verifies at exact rational points of $\R^{8}$ and $\R^{12}$ ($r = 2, 3$: $L^{7}(4)$ and $L^{11}(4)$), for $\epsilon = \pm 1$ and $t \in \{0, \tfrac13, 2\}$: $B$ is orthogonal with $B^{2} = -1$, $f \perp q$, $|f| = |q|$, $f(Bq) = -Bf(q)$, the pull-back of $G_t$ at $Bq$ by $B$ equals $G_t$ at $q$, $G_t(f, f) = t|f|^{2}$, the normal direction $q$ decouples, and the signatures are those stated (on $q^{\perp} \oplus \R$, one positive direction less than on $\R^{4r} \oplus \R$).

\emph{(4).} By (2) and (3), for every odd $n \geq 3$ there is a tame wall with cross-section $L^{2n+1}(4)$, $\operatorname{ht}(x) = 2$ and an indefinite physical side, with $k(\Sig) \geq 1 + \lfloor n/2 \rfloor$, and $k(\Sig) \geq 3$ already at $n = 3$; so $k(\Sig) = \operatorname{ht}(x)$ fails on these walls, $L^{4r-1}(4) \times \R$ with $r \geq 2$ and indefinite side $(4r-1, 1)$, and $k(\Sig) - \operatorname{ht}(x)$ is unbounded. On these walls $\beta x$ is the element of order two of $H^{2}(\Sig; \Z) = \Z_{4}$ (nonzero since $H^{1}(\Sig; \Z) = 0$ and $x \neq 0$), and $(\beta x)^{2} = 0$ always when $x^{2} = 0$: $\beta x$ reduces to $x^{2} = 0$, so $\beta x = 2z$ for an integral $z$, and $2\beta x = 0$ gives $4z = 0$, hence $(\beta x)^{2} = 4z^{2} = 0$. The obstruction of (1) is a statement about stable classes, so it applies in the stable range as well: the failure of $k = \operatorname{ht}(x)$ here does not come from the step of the lens-space proof that needed the stable range. By the parent claim, $k(\Sig) \leq e(\shL) \leq 2n + 2$ on these walls, and on a definite side $k(\Sig) = e(\shL)$: with $q = 0$, $N(\overline G) = 0$ and clause (1) of the parent gives $\shL \oplus N \cong \underline{\R}^{k}$, so $\shL \hookrightarrow \underline{\R}^{k}$ and $k \geq e(\shL)$, while clause (2) gives $k \leq e(\shL)$. The exact values are not determined here. The indefinite side built in (3) has one negative direction, $(2n+1, 1)$, not the physical type $(1, q)$; the obstruction (1) holds on every side, so it applies equally to walls of type $(1, 2n+1)$ wherever they exist.
